% ============================================================
% Appendix C: Lab 2 — Teacher Reference
% ============================================================

\chapter{Lab~2: SYNAPSE Intelligence Portal --- Teacher Reference}
\label{app:lab2}

This appendix provides the deployment reference for Lab~2 (Web Application
Vulnerabilities). It is intended for teachers and administrators. The
application source code, Docker Compose files, and interactive walkthrough
are at \path{docs/web/docs/labs/lab2/}.

\section*{Credentials and Access}

\begin{table}[H]
\centering
\begin{tabularx}{\textwidth}{llXl}
\toprule
\textbf{System} & \textbf{Username} & \textbf{Password} & \textbf{Role} \\
\midrule
pfSense-LAB2         & admin     & pfsense              & Firewall management \\
VyOS-LAB2            & vyos      & vyos                 & Router console \\
Server-A OS          & ubuntu    & \texttt{S3rv3rA}     & SSH access \\
Server-A Portal      & admin     & \texttt{Admin@Synapse2024} & Web application admin \\
Server-A Portal      & guest     & guest123             & Web application guest \\
Server-B OS          & ubuntu    & \texttt{S3rv3rB}     & SSH access \\
Server-B Monitor     & monitor   & \texttt{M0nit0r2024} & Node Monitor API access \\
Parrot-LAB2          & lab2      & L4b2                 & Attacker node \\
\bottomrule
\end{tabularx}
\caption{Lab~2 credentials (teacher reference --- do not distribute to students)}
\label{tab:lab2-creds}
\end{table}

\section*{Network Segments}

\begin{table}[H]
\centering
\begin{tabularx}{\textwidth}{lllX}
\toprule
\textbf{Segment} & \textbf{Subnet} & \textbf{Gateway} & \textbf{Purpose} \\
\midrule
WAN-sim       & 10.0.2.0/24   & 10.0.2.1  & Parrot (10.0.2.10) $\leftrightarrow$ pfSense WAN (10.0.2.1) \\
Net-Link      & 192.168.2.0/24 & ---        & pfSense LAN (192.168.2.1) $\leftrightarrow$ VyOS eth0 (192.168.2.2) \\
DMZ           & 10.0.3.0/24   & 10.0.3.1  & Server-A (10.0.3.10), Server-B (10.0.3.20) via VyOS eth1 \\
\bottomrule
\end{tabularx}
\caption{Lab~2 network segments}
\label{tab:lab2-segments}
\end{table}

\section*{Vulnerability Chain and Flag Locations}

\begin{table}[H]
\centering
\small
\begin{tabularx}{\textwidth}{lllXl}
\toprule
\textbf{Step} & \textbf{OWASP} & \textbf{Vuln} & \textbf{Implementation} & \textbf{Flag} \\
\midrule
1 & A03:2021 & Stored XSS    & \texttt{|safe} filter in Jinja2; \texttt{/comments} endpoint; victim bot auto-visits & --- \\
2 & A07:2021 & Broken auth   & Cookie format \texttt{ID:ROLE:USERNAME}; unsigned; forgeable & --- \\
3 & A03:2021 & SQL injection & \texttt{/search?q=} uses string concatenation; \texttt{admin\_users} table & \texttt{FLAG\{synapse\_sqli\_creds\_dumped\}} \\
4 & A01:2021 & Lateral move  & SSH to Server-B with \texttt{monitor/M0nit0r2024} extracted via SQLi & \texttt{FLAG\{synapse\_pivot\_server\_b\}} \\
5 & A03:2021 & Cmd injection & Server-B \texttt{/ping} uses \texttt{os.popen("ping -c 2 " + host)} & \texttt{FLAG\{synapse\_rce\_achieved\}} \\
\bottomrule
\end{tabularx}
\caption{Lab~2 vulnerability chain and flag locations}
\label{tab:lab2-vuln-chain}
\end{table}

\section*{Lab Startup Procedure}

\begin{enumerate}
    \item After EVE-NG host reboot, restore bridge addresses:
\begin{lstlisting}[language=bash]
bash /usr/local/bin/lab2-bridges-up.sh
\end{lstlisting}
    \item Start the SYNAPSE Portal on Server-A:
\begin{lstlisting}[language=bash]
sshpass -p S3rv3rA ssh ubuntu@10.0.3.10
cd /opt/synapse && docker-compose up -d
docker-compose ps   # expect: flask, nginx, victim -- all Up
\end{lstlisting}
    \item Start the Node Monitor on Server-B:
\begin{lstlisting}[language=bash]
sshpass -p S3rv3rB ssh ubuntu@10.0.3.20
cd /opt/synapse-monitor && sudo docker compose up -d
sudo docker compose ps   # expect: monitor-flask Up on :80
\end{lstlisting}
    \item Configure Parrot static IP and verify connectivity:
\begin{lstlisting}[language=bash]
# On Parrot (10.0.2.10/24, GW 10.0.2.1):
nmcli con mod "Wired connection 1" ipv4.method manual \
  ipv4.addresses 10.0.2.10/24 ipv4.gateway 10.0.2.1
nmcli con up "Wired connection 1"
curl http://10.0.3.10   # Should return SYNAPSE Portal HTML
\end{lstlisting}
    \item Verify all flags are in place:
\begin{lstlisting}[language=bash]
# Flag 1 -- SQLi result (in admin_users table on Server-A)
sshpass -p S3rv3rA ssh ubuntu@10.0.3.10 \
  "sqlite3 /opt/synapse/db/synapse.db 'SELECT flag FROM admin_users'"
# Expected: FLAG{synapse_sqli_creds_dumped}

# Flag 2 -- pivot marker on Server-B
sshpass -p S3rv3rB ssh ubuntu@10.0.3.20 cat /home/monitor/flag.txt
# Expected: FLAG{synapse_pivot_server_b}

# Flag 3 -- RCE result on Server-B
sshpass -p S3rv3rB ssh ubuntu@10.0.3.20 cat /root/flag_final.txt
# Expected: FLAG{synapse_rce_achieved}
\end{lstlisting}
\end{enumerate}

\noindent\textbf{Note:} pfSense-LAB2 is installed in UEFI mode. If it fails to boot,
use the start script: \texttt{/usr/local/bin/pfsense-lab2-start.sh}. Server-A
uses \texttt{docker-compose} (v1, hyphen); Server-B uses \texttt{docker compose}
(v2, space).
